\section{Why radical towers have solvable Galois groups}
\label{sec:radical-solvable}

This section proves the algebraic implication used by Abel--Ruffini.  The proof
is arranged so that every set under discussion is finite and every map can be
recorded by a finite table.  The usual field-theoretic language is retained
for readability; Section~\ref{sec:internalization} replaces each object by its
natural-number code.

\subsection{Finite-group solvability}

For a finite group $G$, let $G'=[G,G]$ be the subgroup generated by all
commutators
\[
 [g,h]=ghg^{-1}h^{-1}.
\]
Define the derived series by
\[
 G^{(0)}=G,
 \qquad
 G^{(i+1)}=(G^{(i)})'.
\]
The group is \emph{solvable} if $G^{(r)}=1$ for some $r$.
For a finite table this is decidable: compute commutators, close under
multiplication and inverse, and iterate.  The subgroup chain can strictly
decrease at most $|G|-1$ times, so $|G|$ iterations suffice.

We use three closure facts, each admitting a finite constructive proof.

\begin{lemma}[Subgroups]
A subgroup of a solvable finite group is solvable.
\end{lemma}

\begin{proof}
If $H\le G$, then $H'\le G'$ because every commutator of elements of $H$ is a
commutator in $G$.  Primitive recursion gives $H^{(i)}\le G^{(i)}$ for all
$i$.  Once $G^{(r)}=1$, also $H^{(r)}=1$.
\end{proof}

\begin{lemma}[Quotients]
A quotient of a solvable finite group is solvable.
\end{lemma}

\begin{proof}
For a surjective homomorphism $\pi:G\to Q$, one has
$\pi(G')=Q'$: images of commutators are commutators, and every commutator in
$Q$ has lifted arguments.  Hence $\pi(G^{(i)})=Q^{(i)}$ by finite induction.
If $G^{(r)}=1$, then $Q^{(r)}=1$.
\end{proof}

\begin{lemma}[Extensions]
Suppose
\[
 1\longrightarrow N\longrightarrow G\longrightarrow Q\longrightarrow1
\]
is exact and both $N$ and $Q$ are solvable.  Then $G$ is solvable.
\end{lemma}

\begin{proof}
If $Q^{(r)}=1$, the quotient calculation gives $G^{(r)}\le N$.  If
$N^{(s)}=1$, then
\[
 G^{(r+s)}=(G^{(r)})^{(s)}\le N^{(s)}=1.
\]
All subgroup inclusions can be verified from finite multiplication tables.
\end{proof}

\subsection{One normal radical layer}

Let $F$ be a characteristic-zero field containing all $m$th roots of unity,
and suppose $E/F$ is a finite normal extension generated by roots of finitely
many polynomials
\[
 Z^m-a_1,\ldots,Z^m-a_k,
 \qquad a_j\in F.
\]
Choose roots $\beta_j$ with $\beta_j^m=a_j$.  If $a_j\ne0$, every root of
$Z^m-a_j$ is $\zeta\beta_j$ for some $m$th root of unity $\zeta\in F$.
If $a_j=0$, the only root is zero and that generator is redundant.

For $\sigma\in\Gal(E/F)$, each nonzero $\beta_j$ is therefore sent to
\[
 \sigma(\beta_j)=\zeta_j(\sigma)\beta_j,
 \qquad \zeta_j(\sigma)^m=1.
\]
The map
\[
 \sigma\longmapsto
 (\zeta_1(\sigma),\ldots,\zeta_k(\sigma))
\]
is an injective group homomorphism into a finite product of cyclic groups of
$m$th roots of unity.  It is injective because the $\beta_j$ generate $E$ over
$F$.  Hence $\Gal(E/F)$ is abelian.

This argument is finite.  A certificate lists the generators, their power
equations, and the root-of-unity multiplier by which each automorphism sends
each generator.  Homomorphism and injectivity are bounded table checks.

\subsection{A normal closure built one radical layer at a time}

Consider a radical tower
\[
 K=K_0\subseteq K_1\subseteq\cdots\subseteq K_r=M,
 \qquad
 K_i=K_{i-1}(\alpha_i),
 \qquad
 \alpha_i^{m_i}=a_i\in K_{i-1}.
 \tag{R1}\label{eq:radical-tower}
\]
Let $m$ be a common multiple of all $m_i$ and put
\[
 N_0=K(\mu_m),
\]
where $\mu_m$ denotes the $m$th roots of unity.  The extension $N_0/K$ is
finite Galois.  An automorphism is determined by the image of a primitive
$m$th root, so its Galois group embeds in the abelian group
$(\Z/m\Z)^\times$ and is therefore abelian.

We now construct finite extensions
\[
 N_0\subseteq N_1\subseteq\cdots\subseteq N_r=N
\]
such that every $N_i/K$ is normal and $N_i$ contains the images of
$N_0K_i$ under all $K$-embeddings into a fixed finite splitting presentation.
Assume $N_{i-1}$ has been constructed.  Let $A_i$ be the finite set
\[
 A_i=\bigl\{\tau(a_i):
   \tau:N_0K_{i-1}\hookrightarrow N_{i-1}
   \text{ is a $K$-embedding}\bigr\}.
\]
For each $b\in A_i$, adjoin one root $\beta_b$ of $Z^{m_i}-b$ and put
\[
 N_i=N_{i-1}(\beta_b:b\in A_i).
\]
Because $N_{i-1}$ contains all $m_i$th roots of unity, every polynomial
$Z^{m_i}-b$ splits in $N_i$.  The set $A_i$ is stable under
$K$-automorphisms of $N_{i-1}$: composing a listed embedding with such an
automorphism merely produces another listed embedding.  Consequently every
$K$-embedding of $N_i$ maps each adjoined root to another root already in
$N_i$.  Thus $N_i/K$ is normal.

The restriction kernel
\[
 H_i=\Gal(N_i/N_{i-1})
\]
is abelian by the preceding subsection: $N_i$ is generated over $N_{i-1}$ by
roots of the finitely many binomials $Z^{m_i}-b$, and $N_{i-1}$ contains the
required roots of unity.  Since both $N_i/K$ and $N_{i-1}/K$ are finite
normal separable extensions, restriction gives an exact sequence
\[
 1\longrightarrow H_i
 \longrightarrow \Gal(N_i/K)
 \longrightarrow \Gal(N_{i-1}/K)
 \longrightarrow1.
 \tag{R2}\label{eq:radical-exact}
\]
Surjectivity is the finite extension-of-embeddings theorem in this
special situation.  Regard an automorphism of $N_{i-1}$ as a $K$-embedding
into $N_i$ and extend it across the finitely generated algebraic extension
$N_i/N_{i-1}$.  Because $N_i/K$ is normal, the extended embedding has image
$N_i$ and is therefore an automorphism.  One must not choose the images of the
$\beta_b$ independently, since they may satisfy relations.  In the coded
proof, a single coordinate map for the extension is supplied; the checker
verifies all generator relations, multiplicativity, and invertibility.  Thus
no choice principle is invoked by the verifier.

Starting with the abelian group $\Gal(N_0/K)$ and repeatedly applying the
extension lemma to~\eqref{eq:radical-exact}, we obtain that
$\Gal(N/K)$ is solvable.  The field $N$ contains the original tower $M$ via
the chosen embedding of each $\alpha_i$ into the corresponding root layer.

\begin{theorem}[Radical-tower theorem]
\label{thm:radical-group-solvable}
Let $M/K$ be a finite extension contained in a finite radical tower of
characteristic-zero fields.  There is a finite normal extension $N/K$
containing $M$ such that $\Gal(N/K)$ is solvable.
\end{theorem}

\begin{proof}
Remove any trivial zero radical generators, adjoin the finitely many required
roots of unity, and apply the layer construction above.  Characteristic zero
gives separability.  The exact sequences have abelian kernels and the initial
cyclotomic group is abelian, so finite induction proves solvability.
\end{proof}

\subsection{Passing to a splitting field}

\begin{corollary}[A polynomial split by radicals has solvable Galois group]
\label{cor:split-rad-solvable}
Let $f\in K[T]$ be separable, and let $L/K$ be its splitting field.  If a
radical tower contains all roots of $f$, then $\Gal(L/K)$ is solvable.
\end{corollary}

\begin{proof}
Take the solvable normal extension $N/K$ from
Theorem~\ref{thm:radical-group-solvable}; it contains $L$.  Since $L/K$ is
normal, restriction gives a surjection
\[
 \Gal(N/K)\twoheadrightarrow\Gal(L/K).
\]
The target is a quotient of a solvable group and is therefore solvable.
\end{proof}

The extension steps are constructive in the certificate sense used here: all
extensions are finite, every relevant polynomial has a finite explicitly
listed root set, and every embedding is extended one generator at a time by a
supplied finite table.  The checker uses neither a maximal-ideal theorem nor a
global algebraic closure.

\subsection{Primitive-recursive bounds}

Let the tower length be $r$ and let the exponents be
$m_1,\ldots,m_r$.  A rough degree bound for the unnormalized tower is
\[
 d=\prod_{i=1}^{r}m_i.
\]
A normal closure and every automorphism group can be bounded by a recurrence
using products, powers, and factorials.  For example, after replacing the
root-of-unity degree by an explicit bound $B_0$, one may use
\[
 B_{i+1}=B_i\,m_i^{B_i}
\]
and finally $B_r!$ as a permutation-table bound.  These estimates are
intentionally extravagant but primitive recursive.  They bound every
enumeration of embeddings, roots, automorphisms, subgroup closures, and
derived-series stages.  Efficiency is irrelevant to the foundational claim.
